% GENERATED FILE. Edit canonical lessons, then run npm run generate:books.
% canonical-source-hash: fnv1a64:cf259a689696fdd7
% canonical-lessons: MW-C277-bhajan, MW-C277-rag, MW-C277-lay, MW-C277-dolno, MW-C277-rijhno

\chapter{Songs and Dancing}
\label{ch:mw-songs-and-dancing}
\hlchaptermodality{277}

\begin{chapteropening}
\textbf{By the end of this chapter:} \emph{I can say which songs and dances I enjoy.}

Complete the last lesson of chapter 277: I can say which songs and dances I enjoy.
\end{chapteropening}

\section[\texorpdfstring{\mw{भजन} (bhajan)}{भजन (bhajan)}]{\mw{भजन} (bhajan) — a devotional song}

\label{lesson:MW-C277-bhajan}



\begin{quote}
Before the new one: say the Marwadi for a clap, then the Marwadi for to ring out.
\end{quote}

\subsection*{You'll want to know: \mw{भजन}}
\textbf{\mw{भजन}} — \emph{bhajan} — \textquotedblleft{}a devotional song\textquotedblright{}.

\begin{quote}
\mw{दादी} \mw{रोज} \mw{भजन} \mw{गावै} \mw{है।}
\end{quote}

\begin{quote}
\emph{dādī roj bhajan gāvai hai.}
\end{quote}

\begin{quote}
Grandmother sings a hymn every day.
\end{quote}

A masculine noun.

\begin{cousinweb}
\textbf{\mw{भजन}} (\emph{bhajan}) is a tatsama, taken from Sanskrit \emph{bhajana}, \textquotedblleft{}reverence, worship\textquotedblright{} (Monier-Williams); in Marwadi it is the song of worship.
\end{cousinweb}

\begin{grammarlens}[title={What you've built}]
One more word for this chapter.
\end{grammarlens}

\subsection*{Your turn}
\begin{itemize}
\raggedright
  \item \emph{Say it:} \emph{bhajan}
  \item \emph{Say it:} \emph{bhajan}, once more
  \item \emph{Say it:} \emph{dādī roj bhajan gāvai hai.} — Grandmother sings a hymn every day.
  \item \emph{Recall:} say the Marwadi for a clap, then the Marwadi for to ring out, then say \emph{bhajan} again
\end{itemize}

\begin{tcolorbox}[breakable,colback=teal!4,colframe=teal!35!black,title={Before you move on}]
What does \mw{भजन} mean? (\textquotedblleft{}A devotional song\textquotedblright{}.) Say it once more.
\end{tcolorbox}
\section[\texorpdfstring{\mw{राग} (rāg)}{राग (rāg)}]{\mw{राग} (rāg) — a melody}

\label{lesson:MW-C277-rag}



\begin{quote}
Before the new one: say the Marwadi for to ring out, then the Marwadi for a devotional song.
\end{quote}

\subsection*{You'll want to know: \mw{राग}}
\textbf{\mw{राग}} — \emph{rāg} — \textquotedblleft{}a melody\textquotedblright{}.

\begin{quote}
\mw{यो} \mw{राग} \mw{जूनो} \mw{है।}
\end{quote}

\begin{quote}
\emph{yo rāg jūno hai.}
\end{quote}

\begin{quote}
This melody is old.
\end{quote}

A masculine noun.

\begin{cousinweb}
\textbf{\mw{राग}} (\emph{rāg}) is a tatsama, taken from Sanskrit \emph{rāga}, which Monier-Williams gives as \textquotedblleft{}a musical note, harmony, melody\textquotedblright{} (Monier-Williams).
\end{cousinweb}

\begin{grammarlens}[title={What you've built}]
One more word for this chapter.
\end{grammarlens}

\subsection*{Your turn}
\begin{itemize}
\raggedright
  \item \emph{Say it:} \emph{rāg}
  \item \emph{Say it:} \emph{rāg}, once more
  \item \emph{Say it:} \emph{yo rāg jūno hai.} — This melody is old.
  \item \emph{Recall:} say the Marwadi for to ring out, then the Marwadi for a devotional song, then say \emph{rāg} again
\end{itemize}

\begin{tcolorbox}[breakable,colback=teal!4,colframe=teal!35!black,title={Before you move on}]
What does \mw{राग} mean? (\textquotedblleft{}A melody\textquotedblright{}.) Say it once more.
\end{tcolorbox}
\section[\texorpdfstring{\mw{लय} (lay)}{लय (lay)}]{\mw{लय} (lay) — rhythm}

\label{lesson:MW-C277-lay}



\begin{quote}
Before the new one: say the Marwadi for a devotional song, then the Marwadi for a melody.
\end{quote}

\subsection*{You'll want to know: \mw{लय}}
\textbf{\mw{लय}} — \emph{lay} — \textquotedblleft{}rhythm\textquotedblright{}.

\begin{quote}
\mw{लय} \mw{तेज} \mw{है।}
\end{quote}

\begin{quote}
\emph{lay tej hai.}
\end{quote}

\begin{quote}
The rhythm is fast.
\end{quote}

A feminine noun.

\begin{cousinweb}
\textbf{\mw{लय}} (\emph{lay}) is a tatsama, taken from Sanskrit \emph{laya}, which Monier-Williams gives, in music, as \textquotedblleft{}time\textquotedblright{} (Monier-Williams).
\end{cousinweb}

\begin{grammarlens}[title={What you've built}]
One more word for this chapter.
\end{grammarlens}

\subsection*{Your turn}
\begin{itemize}
\raggedright
  \item \emph{Say it:} \emph{lay}
  \item \emph{Say it:} \emph{lay}, once more
  \item \emph{Say it:} \emph{lay tej hai.} — The rhythm is fast.
  \item \emph{Recall:} say the Marwadi for a devotional song, then the Marwadi for a melody, then say \emph{lay} again
\end{itemize}

\begin{tcolorbox}[breakable,colback=teal!4,colframe=teal!35!black,title={Before you move on}]
What does \mw{लय} mean? (\textquotedblleft{}Rhythm\textquotedblright{}.) Say it once more.
\end{tcolorbox}
\section[\texorpdfstring{\mw{डोलणो} (ḍolṇo)}{डोलणो (ḍolṇo)}]{\mw{डोलणो} (ḍolṇo) — to sway; to wander}

\label{lesson:MW-C277-dolno}



\begin{quote}
Before the new one: say the Marwadi for a melody, then the Marwadi for rhythm.
\end{quote}

\subsection*{You'll want to know: \mw{डोलणो}}
\textbf{\mw{डोलणो}} — \emph{ḍolṇo} — \textquotedblleft{}to sway; to wander\textquotedblright{}.

\begin{quote}
\mw{छोरी} \mw{नाचै} \mw{है} \mw{अर} \mw{डोलै} \mw{है।}
\end{quote}

\begin{quote}
\emph{chhorī nāchai hai ar ḍolai hai.}
\end{quote}

\begin{quote}
The girl dances and sways.
\end{quote}

Past (he …): \mw{वो} \mw{गांव} \mw{मांय} \mw{डोल्यो}, \emph{vo gā̃v mā̃y ḍolyo}.

\begin{cousinweb}
\textbf{\mw{डोलणो}} (\emph{ḍolṇo}) is inherited from Sanskrit \emph{dolāyate}, \textquotedblleft{}moves to and fro\textquotedblright{}; CDIAL lists Old Marwari \emph{ḍolaï}, \textquotedblleft{}wanders\textquotedblright{} (CDIAL 6585).
\end{cousinweb}

\begin{grammarlens}[title={What you've built}]
One more word for this chapter.
\end{grammarlens}

\subsection*{Your turn}
\begin{itemize}
\raggedright
  \item \emph{Say it:} \emph{ḍolṇo}
  \item \emph{Say it:} \emph{ḍolṇo}, once more
  \item \emph{Say it:} \emph{chhorī nāchai hai ar ḍolai hai.} — The girl dances and sways.
  \item \emph{Recall:} say the Marwadi for a melody, then the Marwadi for rhythm, then say \emph{ḍolṇo} again
\end{itemize}

\begin{tcolorbox}[breakable,colback=teal!4,colframe=teal!35!black,title={Before you move on}]
What does \mw{डोलणो} mean? (\textquotedblleft{}To sway; to wander\textquotedblright{}.) Say it once more.
\end{tcolorbox}
\section[\texorpdfstring{\mw{रीझणो} (rījhṇo)}{रीझणो (rījhṇo)}]{\mw{रीझणो} (rījhṇo) — to be delighted}

\label{lesson:MW-C277-rijhno}



\begin{quote}
Before the new one: say the Marwadi for rhythm, then the Marwadi for to sway.
\end{quote}

\subsection*{You'll want to know: \mw{रीझणो}}
\textbf{\mw{रीझणो}} — \emph{rījhṇo} — \textquotedblleft{}to be delighted\textquotedblright{}.

\begin{quote}
\mw{दादी} \mw{गीत} \mw{सुणै} \mw{है} \mw{अर} \mw{रीझै} \mw{है।}
\end{quote}

\begin{quote}
\emph{dādī gīt suṇai hai ar rījhai hai.}
\end{quote}

\begin{quote}
Grandmother listens to the song and is delighted.
\end{quote}

Past (he …): \mw{वो} \mw{रीझ्यो}, \emph{vo rījhyo}.

\begin{cousinweb}
\textbf{\mw{रीझणो}} (\emph{rījhṇo}) is inherited from Sanskrit \emph{ṛdhyati}, \textquotedblleft{}increases, prospers\textquotedblright{}; CDIAL lists Old Marwari \emph{rījhaï}, \textquotedblleft{}is satisfied\textquotedblright{}, and Hindi \emph{rījhnā}, \textquotedblleft{}to be enchanted\textquotedblright{} (CDIAL 2457).
\end{cousinweb}

\begin{grammarlens}[title={What you've built}]
One more word for this chapter.
\end{grammarlens}

\subsection*{Your turn}
\begin{itemize}
\raggedright
  \item \emph{Say it:} \emph{rījhṇo}
  \item \emph{Say it:} \emph{rījhṇo}, once more
  \item \emph{Say it:} \emph{dādī gīt suṇai hai ar rījhai hai.} — Grandmother listens to the song and is delighted.
  \item \emph{Recall:} say the Marwadi for rhythm, then the Marwadi for to sway, then say \emph{rījhṇo} again
\end{itemize}

\begin{tcolorbox}[breakable,colback=teal!4,colframe=teal!35!black,title={Before you move on}]
What does \mw{रीझणो} mean? (\textquotedblleft{}To be delighted\textquotedblright{}.) Say it once more.
\end{tcolorbox}
